\chapter*{RESUMEN}
\addcontentsline{toc}{chapter}{RESUMEN}

El presente perfil de proyecto de grado tiene como propósito fundamentar y estructurar el diseño de un sistema informático de gestión y trazabilidad de servicios de soporte técnico con asistencia mediante inteligencia artificial, orientado a prestadores de servicios informáticos en la ciudad de Tarija durante la gestión 2026. La investigación parte de la necesidad de analizar cómo los prestadores de servicios de soporte técnico de la ciudad de Tarija gestionan la recepción de equipos, el registro de diagnósticos, la comunicación de presupuestos y el seguimiento de las reparaciones. Mediante el trabajo de campo se identificarán las modalidades de registro utilizadas y las posibles dificultades de comunicación y consulta de información por parte de los clientes, con el propósito de fundamentar los requerimientos del sistema propuesto. La metodología adopta un enfoque mixto (cualitativo y cuantitativo), de nivel descriptivo-propositivo y diseño no experimental transversal, aplicando un muestreo no probabilístico por conveniencia sobre dos unidades de análisis: 10 a 15 técnicos o propietarios de talleres y 50 a 100 usuarios. La solución propuesta se adscribe al área de la ingeniería de software bajo la arquitectura Modelo-Vista-Controlador (MVC) con el framework PHP Laravel, MySQL y una interfaz responsiva con Bootstrap 5, donde el diseño contempla un módulo de trazabilidad mediante estados secuenciales, evidencias fotográficas y restricciones de borrado en base de datos. Asimismo, se propone integrar un asistente conversacional basado en inteligencia artificial estructurado sobre una API REST de un proveedor externo, tomando Google Gemini como alternativa tecnológica de referencia, cuya selección definitiva se realizará considerando su disponibilidad, continuidad del servicio, coste de utilización y capacidad para procesar consultas en español. El backend de Laravel validará las solicitudes, consultará únicamente los datos autorizados y proporcionará al asistente el contexto necesario para responder preguntas frecuentes y consultas de estado mediante código verificado como característica prevista en la solución. La propuesta se fundamenta en marcos de gobernanza como ITIL y estándares ISO/IEC 25010 e ISO/IEC 27001.